
% \chapter{Problema 3.3}

% \section{Solução analítica}

% Temos a equação diferencial ordinária (EDO) dada por:

% \[ \od{y}{t} = \left( \alpha + j\omega \right)y \]

% \noindent e o objetivo é identificar a solução $y(t)$.

% Usa separação de variáveis 

% \[ \frac{1}{y} dy  = \left( \alpha + j\omega \right) dt  \]

% Integrando ambos os lados

% \[ \int_{t_0}^{t} \frac{1}{y} dy  = \int_{t_0}^{t} \left( \alpha + j\omega \right) dt  \]

% \[ \ln|y(t)| - \ln|y(t_0)| = \left( \alpha + j\omega \right)(t - t_0) \]

% Considerando $t_0 = 0$ e $y(t_0) = y_0$, temos

% \[ \ln|y(t)| = \ln|y_0| + \left( \alpha + j\omega \right)t \]

% \[ y(t) = y_0 e^{\left( \alpha + j\omega \right)t} \]

% Considerando $y(t_0) = 1$, 

% \[ y(t) = e^{\left( \alpha + j\omega \right)t} \]

% Plotando essa solução, temos 

% \begin{figure}[htb]
% 	\caption{\label{fig:analitica_q1}Solução analítica da questão 1.}
% 	\begin{center}
%         \includegraphics[width=\textwidth,trim=1 1 1 1,clip]{analitica_q1.png}
% 	\end{center}
% \end{figure}

% \section{Solução numérica via Método de Euler}

% Segundo o método de Euler, temos 

% \[ y_{n+1} = y_n + \Delta t \left( \alpha + j\omega \right) y_n \]

% \noindent com $y_0 = 1$. Considerando $\omega = 2$, $\alpha = 0$,
% $\Delta = 0.1$, 

% \[ y_{n+1} = y_n + j \, 2 \, 0.1 \,   y_n \]

% \[ y_{n+1} = y_n + j \, 0.2 \, y_n \]

% Resolvendo numericamente para $t \in [0, 10]$, temos a seguinte solução numérica

% \begin{figure}[htb]
% 	\caption{\label{fig:euler_q1}Solução numérica via Método de Euler da questão 1.}
% 	\begin{center}
%         \includegraphics[width=\textwidth,trim=1 1 1 1,clip]{euler_q1.png}
% 	\end{center}
% \end{figure}


% \section{Solução numérica via Método Leapfrog}


\chapter{Problema 3.4}

\section{Solução analítica}

A solução analítica vem da equação de D'Alembert, conforme descrito 
no livro-text: 

\[ \mathcal{V}(x, t) = \frac{1}{2} \left[ \phi(x - v_p t) + \phi(x + v_p t) + \frac{1}{v_p} \int_{x - v_p t}^{x + v_p t} \theta(\zeta) d\zeta \right] \]

Temos 

\begin{itemize}
	\item $\phi$ é a função da condição inicial de $\mathcal{V}(x, t)$, dada no enunciado
	\item $\theta$ é a função da condição inicial de $\mathcal{V}(x, t)$, que assumimos ser zero
	uma vez que a corrente inicial é nula
\end{itemize}

Assim, 

\[ \mathcal{V}(x, t) = \frac{1}{2} \left[ \phi(x - v_p t) + \phi(x + v_p t) \right] \]  

Temos duas ondas se propagandos em direções opostas com metade da amplitude inicial.

\section{Solução numérica}

O código usado nessa solução encontra-se no Anexo \ref{cap:codigo_q34}.

Os parâmetros de discretização foram selecionados da seguinte forma:

\begin{itemize}
	\item $\Delta t = 5 \un{ps}$: isso vem do fato de que a velocidade da onda é de $v_p = 1.25 \cdot 10^{10} \un{cm/s}$. 
	Como o comprimento da linha é de 20 cm, o tempo de propagação é de $t_p = \frac{20 \un{cm}}{1.25 \cdot 10^{10} \un{cm/s}} = 1.6 \un{ns}$.
	Assim, um passo de tempo de $5 \un{ps}$ é suficientemente pequeno para capturar a dinâmica da propagação da onda.
	\item $\Delta x = 0.1 \un{cm}$: isso é escolhido para garantir que tenhamos uma resolução espacial adequada para capturar a forma da onda, obtendo-se 200 pontos espaciais ao longo da linha de 20 cm.
\end{itemize}

A Figura \ref{fig:result_q34} apresenta a solução numérica para diferentes instantes de tempo.
Ela mostra como a onda se propaga como duas ondas viajando em direções opostas, uma para a direita e outra para a esquerda, a partir da condição inicial dada,
conforme previsto na solução analítica. 

\begin{figure}[htb]
	\caption{\label{fig:result_q34}Solução numérica da questão 3.4.}
	\begin{center}
        \includegraphics[width=\textwidth,trim=1 1 1 1,clip]{result_q34.png}
	\end{center}
\end{figure}

\section{Variando os parâmetros}